\section{A chart at the diagonal}\label{sec:diagonal}

The construction of Section \ref{sec:slices} cannot be applied
directly at $\Delta$: the $\SO(3)$ representation has rank-one
zeros with circle stabiliser. We construct a chart near that
rank-one locus and a second chart away from it, then join them
through the free part. The normalisation near rank one is what
ensures the stronger normal estimate in \eqref{eq:edge-split}.

\subsection{Scaling the slice}
Use Proposition \ref{prop:slice} at $p\in\Delta$ and perform
only the fibrewise linear normalisation of Lemma
\ref{lem:normalise}, which requires no freeness. In
\[
 \EE=\EE_0\oplus(\RR^3\otimes\HH^{k+1}),\qquad y=(u,n),
 \qquad r=|n|,
\]
the resulting data $\mathcal D=(\widehat g,\widehat I_a,
\widehat\mu)$ have flat normal blocks along $\EE_0$, satisfy
\eqref{eq:slice-moment} on their zero set, and obey
\[
 \widehat\mu=\mu_0+O(r^3),\qquad
 \widehat\mu|_{\EE^T}=0
\]
for every maximal circle $T\subset\SO(3)$. The latter follows
from $\mathfrak{so}(3)^T=\mathfrak t$ in Proposition
\ref{prop:slice}. Compare with the product data
\[
 \mathcal D_*=(g_B\oplus g_{\fl},
 I_{a,B}\oplus I_a,\mu_0),
\]
where $(g_B,I_{a,B})$ is the restriction to $\EE_0$.
Their quotient is $X_0=(\EE_0,g_B)\times C_2$.

At a normal radius $\bar r>0$, use coordinates
$(u-u_0,n)/\bar r$ and divide the moment map by $\bar r^2$.
Smoothness, normal agreement, and the cubic remainder imply
\begin{equation}\label{eq:scaled-close}
 \|\mathcal D-\mathcal D_*\|_{C^i}\le C_i\bar r
\end{equation}
on fixed-size scaled sets, for every fixed $i$. This includes
the moment map and its derivatives: every term in its remainder
has at least three normal factors.

\subsection{Separating the circle stabiliser}
Fix a unit vector $e\in\RR^3$, its stabiliser $T$, and
$N(T)\cong O(2)$. Write
$\mathfrak{so}(3)=\RR e\oplus e^\perp=\mathfrak t\oplus\mathfrak p$.
For $n:\RR^3\to\HH^{k+1}$ let $s_1\ge s_2\ge s_3$
be its singular values and put $\rho_T=(s_2^2+s_3^2)^{1/2}$.
This is the Euclidean distance to rank at most one.
For a fixed $c_0<1$, the region $\rho_T<c_0s_1$ is
\begin{equation}\label{eq:svd-slice}
 \SO(3)\times_{N(T)}
 \{(u,l,n''):l\ne0,\ n''\in e^\perp\otimes l^{\perp_{\RR}},
                         \ |n''|<c_0|l|\}.
\end{equation}
Indeed the largest eigenline of the Gram matrix is unique there.
Move it to $\RR e$; the remaining ambiguity is precisely $N(T)$.
In these coordinates $n=e\otimes l+n''$, $s_1=|l|$ and
$\rho_T=|n''|$.

Decompose the last normal space as
\[
 e^\perp\otimes l^{\perp_{\RR}}
 =F_l\oplus(e^\perp\otimes\operatorname{Im}(\HH)l),
 \qquad F_l=e^\perp\otimes(\HH l)^\perp.
\]
The group $T$ acts on $F_l$ as on $T^*\CC^{2k}$. Let
$\mathfrak B=\EE_0\times(\HH^{k+1}\setminus\{0\})$ and let
$\mathfrak F$ be the bundle of these $F_l$ with $|f|<c_0|l|$.
If $f=e_2\otimes f_2+e_3\otimes f_3$, the flat equations
along $\mathfrak p$ impose orthogonality to $I_al$, and the
remaining equations are
\[
 \mu_T(f)_a=g_{\fl}(I_af_2,f_3)=0.
\]
Consequently, with $Y_0=\EE_0\times\mu_0^{-1}(0)$,
\begin{equation}\label{eq:zero-slice}
 (Y_0\cap\{\rho_T<c_0s_1\})/\SO(3)
 =\{(b,f)\in\mathfrak F:\mu_T(f)=0\}/N(T).
\end{equation}
On this part of $X_0$, the distance to the singular set equals
$\rho_T$. The Euclidean distance is a lower bound for the
quotient length, while the path
$(u,e\otimes l+tf)$ lies in $Y_0$ and attains that bound.

\begin{lemma}\label{lem:reduced-data}
For $c_0<1/\sqrt2$ and $r$ sufficiently small, the equations
along $\mathfrak p$ determine a smooth equivariant correction
$c^\flat(b,f)\in e^\perp\otimes\operatorname{Im}(\HH)l$.
The map
$j(b,f)=(u,e\otimes l+f+c^\flat(b,f))$ carries the zero set
of $\mu^\flat=(\widehat\mu\circ j)^{\mathfrak t}$ to $Y$
in the slice \eqref{eq:svd-slice}. Projection off the
quaternionic span of $\mathfrak p j(b,f)$ induces smooth
quaternionic Hermitian tensors $(g^\flat,I_a^\flat)$ on
$\mathfrak F$. On $\{\mu^\flat=0\}$ they satisfy the moment
identity for $T$, and their circle quotient is the original
$\SO(3)$ quotient. Along $\mathfrak B$ their moment map has
the quadratic term given by their fibre structure.
\end{lemma}
\begin{proof}
The derivative in $c^\flat$ of the flat $\mathfrak p$ equations
is the map with components $g(I_al,c_j^\flat)$, $j=2,3$.
It is invertible with inverse of size $C/|l|$. Equation
\eqref{eq:scaled-close} and the implicit function theorem give
$c^\flat$, with scaled $C^i$ norm $O(\bar r)$.
It vanishes at $f=0$ since $\widehat\mu$ does on $\EE^T$.
Every nearby zero in the singular-value slice has this form,
by the same uniqueness. For the product data the correction is
identically zero.

Let $Q_z=\mathfrak pz+\sum_a\widehat I_a\mathfrak pz$
and let $p_z$ project off it. We check that $p_zdj$ is
invertible. In the flat model its possible kernel consists of
tangent vectors of the form
\[
 \dot l=\alpha_2f_3-\alpha_3f_2,\qquad
 \dot f_2=\alpha_3l,\qquad \dot f_3=-\alpha_2l
 \quad(\alpha_2,\alpha_3\in\HH).
\]
Differentiating $f_j\perp\HH l$ gives
\[
 (|\alpha_2|^2+|\alpha_3|^2)|l|^2
 \le2|f|^2(|\alpha_2|^2+|\alpha_3|^2).
\]
Thus the kernel is zero for the stated $c_0$. Dimensions agree:
$Q_z$ has real dimension eight. Homogeneity and compactness
give a uniform inverse at scale $|l|$, and
\eqref{eq:scaled-close} preserves invertibility for the actual
data. Define
\[
 g^\flat(v,w)=\widehat g(p_zdjv,p_zdjw),\qquad
 I_a^\flat=(p_zdj)^{-1}\widehat I_a p_zdj.
\]
They differ from the corresponding model data by
$O(\bar r)$ in scaled $C^i$ norm.

At a zero $z$, the four quaternionic orbit subspaces are
orthogonal. Since the $\mathfrak p$ moment equations vanish
identically on the image of $j$, the component of $djv$ in
$Q_z$ lies in $\mathfrak pz$. It is orthogonal to
$\widehat I_a\mathfrak tz$. Equivariance therefore gives
\[
 \langle d\mu_a^\flat(v),\xi\rangle
 =g^\flat(I_a^\flat\xi(b,f),v)\quad(\xi\in\mathfrak t).
\]
Minimising the metric first over $\mathfrak p$ and then over
$\mathfrak t$ identifies the quotient metric, and the same
horizontal lift identifies the complex structures. No reduction
by a subgroup with Lie algebra $\mathfrak p$ is asserted.

It remains to verify the Hessian along $f=0$. There
$\mu^\flat=d\mu^\flat=0$. Let $Q_b$ be its quadratic term.
It is $O(\bar r)$-close to $\mu_T$ after scaling. Every
regular nonzero zero $f$ of $Q_b$ continues to a curve
$(b,tf(t))$ of zeros of $\mu^\flat$ with $f(0)=f$, by
dividing the equation by $t^2$ and applying the implicit function
theorem. Differentiating the preceding moment identity along
this curve gives
\[
 \partial^2\mu_a^\flat(b)(f,f')
 =g_b^\flat(I_a^\flat\xi f,f').
\]
Choose a finite spanning collection of regular zeros of the flat
circle moment map. Such a collection exists: in $T^*\CC^{2k}$
the vectors $(\alpha z_i,\beta z_j^*)$ with $i\ne j$ and
$|\alpha|=|\beta|=1$ are zeros and span the real vector space.
They persist and still span for $Q_b$. The displayed identity,
linear in $f$, therefore holds for every $f$, proving the
quadratic-term assertion.
\end{proof}

\subsection{The chart near rank one}
Over a ball of radius $c|l_0|$ in $\mathfrak B$, identify
$F_{l_0}$ with $F_l$ by the polar part of orthogonal projection
$(\HH l_0)^\perp\to(\HH l)^\perp$, tensored with $e^\perp$.
Denote this quaternionic trivialisation by $\vartheta$.
It commutes with $T$ and has uniform scaled derivatives.
The fibrewise construction of Lemma \ref{lem:normalise}, now
with parameter $b=(u,l)$, gives a map $\Xi^\flat$ which
fixes $\mathfrak B$, takes $\mu_T^{-1}(0)$ to
$(\mu^\flat)^{-1}(0)$, and normalises the fibre metric and
all three complex structures there. Its deviation from the
identity has scaled $C^{1,1}$ norm $O(\bar r)$.

Here is why that estimate is uniform even at $f=0$. After the
linear fibre isometry, the residual map is
$\mu_T+R^\flat$, with
\[
 |\partial^iR^\flat|\le C_i\bar r|f|^{3-i}
\]
in scaled coordinates. The right inverse of $d\mu_T(f)$ has
size $C/|f|$. The same contraction equation as in Lemma
\ref{lem:normalise} gives a correction with derivatives
$C_i\bar r|f|^{2-i}$. Its second derivative is uniformly
bounded; away from $f=0$ all orders are controlled. Every
operation is canonical for the invariant data, hence equivariant
under $N(T)$.

For the model data this map is the identity. Both its reduced
tensors and the pulled-back actual reduced tensors have the same
normal blocks along $\mathfrak B$. Their difference is
$O(\bar r)$ in scaled Lipschitz norm. The relative part of
Lemma \ref{lem:quotient-estimate} gives
\begin{equation}\label{eq:inner-estimate}
 (\Theta^{\rm in})^*g-g_0=t_b+O(\rho_T),\qquad
 |t_b|\le C\bar r,\quad|\nabla t_b|\le C,
 \quad|\nabla^0((\Theta^{\rm in})^*g-g_0)|\le C,
\end{equation}
and the same assertions for each complex structure. Here
$\Theta^{\rm in}$ is induced by $j\Xi^\flat$, and $t_b$ is
pulled back from the smooth rank-one stratum. To check the
scaling, the relative lemma gives an error
$C\bar r\widehat\rho$ at unit scale. Since
$\widehat\rho=\rho_T/\bar r$, it becomes $C\rho_T$.
A covariant derivative costs $\bar r^{-1}$, cancelling
the factor $\bar r$.

Applying the absolute part of the same lemma to the model
reduced tensors gives uniform edge charts
$B_b\times C_1\to X_0$. The smooth factors may be extended
to compact factors with uniform scaled geometry by cutoffs.
The complex structure is extended as an almost complex structure,
which is all the edge comparison uses. For example, write it
locally as $A I(b_0)A^{-1}$ with $A$ close to the identity and
replace $A$ by $1+\chi(A-1)$. The tensors $t_b$ are extended
with the same cutoff. Thus \eqref{eq:inner-estimate} is exactly
the normal and tangential decomposition needed in
Definition \ref{def:two-level}.

\subsection{Joining the charts}
Choose $c_3>0$ smaller than a quarter of the radius of validity
of the inner construction. On the cone $\rho_T>c_3s_1/2$
the action is free near the flat zero set. The construction of
Lemma \ref{lem:normalise} gives a map $\Xi_2$ with
$\mu_0\Xi_2=\widehat\mu$ and scaled $C^i$ distance
$O(\bar r)$ from the identity. Its inverse $\Xi^{\rm out}$
carries $Y_0$ to $Y$ in this region.

Extend the inner map to ambient slice coordinates by
\[
 (b,f,c')\longmapsto
 (u,e\otimes l+f'+c^\flat(b,f')+c'),
 \qquad (b,f')=\Xi^\flat(b,f),
\]
and then equivariantly using \eqref{eq:svd-slice}. It is
smooth on the regular zero fibre, fixes the rank-one locus, and
has scaled $C^i$ distance $O(\bar r)$ from the identity on the
free overlap. No smoothness of this ambient extension is needed
at points with $f=0$ and $c'\ne0$.
On $c_3s_1\le\rho_T\le2c_3s_1$, set
$\Gamma=\Xi_2\Xi^{\rm in}$. This map preserves $Y_0$.

Choose a cutoff $\varphi$ of $\rho_T/s_1$ equal to one at
the inner end and zero at the outer end. The naive interpolation
$x+\varphi(\Gamma x-x)$ need not preserve the zero set.
Correct it to
\[
 G(x)=x+\varphi(\Gamma x-x)+b(x)
\]
by requiring
\begin{equation}\label{eq:glue-moment}
 \mu_0(G(x))=(1-\varphi)\mu_0(x)+\varphi\mu_0(\Gamma x).
\end{equation}
Take $b(x)$ in the image of the right inverse of $d\mu_0(x)$.
Before correction the error in \eqref{eq:glue-moment} is
$(\varphi^2-\varphi)\mu_0(\Gamma x-x)$, because $\mu_0$
is quadratic. It is $O(\bar r^2)$ in scaled $C^i$ norm.
The right inverse is uniformly bounded at unit scale on this
free overlap. Contraction therefore gives a unique correction
with that bound. It vanishes where $\varphi$ is zero or one.
Consequently the map
\[
 \Xi=\Xi^{\rm in}\quad(\rho_T\le c_3s_1),\qquad
 \Xi=\Xi^{\rm out}G\quad(c_3s_1\le\rho_T\le2c_3s_1),\qquad
 \Xi=\Xi^{\rm out}\quad(\rho_T\ge2c_3s_1)
\]
is well defined and carries $Y_0$ into $Y$.

We verify that it is a quotient chart. In the outer region and
overlap its unscaled estimates are
\[
 |\Xi-\id|\le Cr^2,\qquad |d\Xi-1|\le Cr,
 \qquad |d^2\Xi|\le C.
\]
On each rescaled free annulus it is a local diffeomorphism close
to the identity. Uniform local injectivity and its small
displacement imply injectivity on that region. In the inner
region injectivity follows from the unique singular-value slice
and the fibrewise coordinate changes. The smaller inner region
$\rho_T/s_1<c_3/2$ and the region $\rho_T/s_1>c_3$ have
images separated by $c_3/2-O(r)$ in this ratio; the remaining
overlap is already covered by the free-annulus argument. The
equivariant map thus induces an injective quotient map. It is
open away from the vertex stratum, by the local diffeomorphism
property on the free part and the explicit inner product charts
near rank one.

It is also locally onto. On a small set
$|u|<t_1$, $0<r<t_2$, the estimate $|\Xi-\id|\le Cr^2$
and preservation of $u$ make the map proper over
$|u|<t_1$, $0<r<t_2/2$. Its image there is open and closed.
The target is locally a normal analytic space; a small connected
neighbourhood with its proper analytic diagonal removed is
connected. It meets the image, so is contained in it. At $r=0$
both the map and its inverse extend continuously by the same
quadratic displacement estimate. Stabilisers are preserved:
rank-one points are fixed, and the nonzero circle-normal
coordinate stays nonzero. This proves the homeomorphism and
regular-locus diffeomorphism assertions.

\begin{theorem}\label{thm:diagonal-chart}
At every point of $\Delta$, $M$ admits a chart adapted to both
strata, for all three complex structures, with model
$(\EE_0,g_B)\times C_2$.
\end{theorem}
\begin{proof}
Compose the quotient map just constructed with the initial fibre
normalisation and the gauge slice map. In the outer region
Lemma \ref{lem:quotient-estimate} gives $O(r)$ error and
bounded first derivative. In the inner region
\eqref{eq:inner-estimate} gives these bounds and
\eqref{eq:edge-split}. An edge chart may also meet the outer
region; there $\rho\ge c_3s_1\asymp r$, so the same smooth
tensor $t_b$ leaves an error $O(r)=O(\rho)$. Thus the
decomposition holds throughout each edge chart. Shrinking $r$
makes the metrics uniformly equivalent and $|t_b|<1/2$.
The compact extension of the smooth vertex factor is available
for each K\"ahler structure. All conditions of Definition
\ref{def:two-level} follow.
\end{proof}

\begin{corollary}\label{thm:T2}
The Gaffney domain of $(M_{\reg},g)$ is Weil invariant for
each $I_a$. Its harmonic forms carry the pointwise
$\mathfrak{so}(4,1)$ action.
\end{corollary}
\begin{proof}
Proposition \ref{prop:models} gives Weil invariance away from
$\Delta$ and on the models $B\times C_2$. Theorem
\ref{thm:two-level} transfers it through Theorem
\ref{thm:diagonal-chart}. Locality covers $M$. The Stokes
property was proved in Proposition \ref{prop:models}; apply
Proposition \ref{prop:global} to the three structures and use
\eqref{eq:verbitsky}.
\end{proof}
